\documentclass[11pt]{article}
\usepackage[a4paper,margin=2.4cm]{geometry}
\usepackage{amsmath,amssymb}
\usepackage{slashed}
\usepackage{graphicx}
\usepackage{booktabs}
\usepackage{array}
\usepackage{longtable}
\usepackage{xcolor}
\usepackage[hidelinks]{hyperref}
\usepackage{siunitx}
\usepackage{caption}
\setlength{\parindent}{0pt}
\setlength{\parskip}{4pt}

\newcommand{\Zp}{Z^{\prime}}
\newcommand{\gD}{g_{D}}
\newcommand{\mZp}{m_{\Zp}}
\newcommand{\Treh}{T_{\mathrm{reh}}}
\newcommand{\betaH}{\beta/H}
\newcommand{\OmGW}{h^{2}\Omega_{\mathrm{GW}}}

\title{\textbf{Can a classically scale-invariant dark $U(1)$ model explain the\\
NANOGrav signal through a cosmological first-order phase transition?}\\[4pt]
\large A pipeline report on the \texttt{scale-invariant-dark-u1-zprime} model (\texttt{fopt-pta} branch)}
\author{DarkAgents}
\date{2026-06-10}

\begin{document}
\maketitle

\begin{abstract}
We report an end-to-end automated analysis of a classically scale-invariant
(Coleman--Weinberg) gauged dark $U(1)_D$ model with a complex dark scalar $\Phi$,
a dark gauge boson $\Zp$, and a vector-like pair of Dirac dark fermions
($\chi$ and a mirror partner $\eta$), as a candidate explanation of the
NANOGrav 15-year (NG15) nano-hertz stochastic gravitational-wave background (SGWB)
via a strongly supercooled first-order phase transition (FOPT).
The semi-analytic FOPT backend yields a large viable region; the PTA stage finds
posteriors that are \emph{interior} to the (user-enlarged) priors, with a maximum
a posteriori (MAP) point scoring $14/14$ free-spectrum violin bins and parameters
($\gD\simeq0.59$, $y\simeq0.38$, $w\simeq1$~GeV) consistent with the closest
literature precedent (arXiv:2502.19478). We conclude that the model
\emph{can fit} the NG15 free spectrum in the chosen $w\in[0.05,10]$~GeV window.
However, the cosmological viability of that fit is \emph{conditional and unverified}:
the analysis fixed the kinetic-mixing portal $\epsilon=0$ (secluded dark sector),
whereas Standard-Model (SM) reheating and BBN require $\epsilon\sim10^{-6}\text{--}10^{-3}$;
a factor-$\sim7$ $\betaH$ convention mismatch and a saturated GW template ($K_{sw}\to1$)
weaken the spectral-amplitude interpretation; and no Bayes factor was computed.
All five most-decisive cosmological constraints require backend modules that were not run.
\end{abstract}

\section{Prompt}
The original user request, verbatim:
\begin{quote}
\itshape ``Can a classically scale-invariant $U(1)$ model with a dark scalar, a dark gauge
boson and a dark fermion can explain the NANOGrav signal through a cosmological FOPT?''
\end{quote}
During the run the user approved each stage and explicitly requested enlarging the
PTArcade prior upper bounds to $\gD\le2.0$, $y\le2.0$, $w\le10$~GeV
(with the hard cut $y^{4}<1.5\,\gD^{4}$). We report below that the resulting posteriors
are interior and the enlarged edges carry no posterior mass.

\section{Model description}
\label{sec:model}
The model is a classically scale-invariant (no tree-level mass scales; only dimension-4
operators) gauged dark abelian symmetry $U(1)_D$ with coupling $\gD$, under which the SM
is a singlet. Symmetry breaking is radiative (Coleman--Weinberg, CW). The global
$U(1)_\chi$ dark-fermion number is conserved, stabilising the lightest dark fermion.

\paragraph{Field content (Table~\ref{tab:fields}).}
A complex scalar $\Phi$ ($q=1$), whose radial mode is the pseudo-dilaton (scalon) and
whose angular mode is eaten by $\Zp$; the dark gauge boson $\Zp$; and two Weyl pairs
$(\psi_L,\psi_R)$ and $(\eta_L,\eta_R)$ forming two degenerate Dirac fermions
($\chi,\eta$). The mirror $\eta$ is introduced purely so that the chiral
$U(1)_D$ charges render the sector vector-like and exactly anomaly-free.

\begin{table}[h]
\centering
\small
\caption{Field content and $U(1)_D$ charges (from \texttt{handoff\_model.json}).}
\label{tab:fields}
\begin{tabular}{lcccc}
\toprule
Field & Spin & $U(1)_D$ charge & dof & Role \\
\midrule
$\Phi$    & 0   & $+1$  & 2 & complex scalar; radial = scalon, angular eaten \\
$\Zp$     & 1   & 0     & 3 & dark gauge boson, $\mZp=\gD\,w$ \\
$\psi_L$  & 1/2 & $+1/2$ & 2 & Dirac $\chi$ (left) \\
$\psi_R$  & 1/2 & $-1/2$ & 2 & Dirac $\chi$ (right) \\
$\eta_L$  & 1/2 & $-1/2$ & 2 & mirror (anomaly partner) \\
$\eta_R$  & 1/2 & $+1/2$ & 2 & mirror (anomaly partner) \\
\bottomrule
\end{tabular}
\end{table}

\paragraph{Lagrangian.}
With $D_\mu\Phi=(\partial_\mu-i\gD \Zp_\mu)\Phi$ and $\Phi=\tfrac{1}{\sqrt2}\chi\,e^{i\theta}$,
\begin{align}
\mathcal{L} =\;& -\tfrac14 F^{\prime}_{\mu\nu}F^{\prime\,\mu\nu}
 + |D_\mu\Phi|^2
 + i\,\bar\psi\,\slashed{D}\,\psi + i\,\bar\eta\,\slashed{D}\,\eta \nonumber\\
 & -\big(y\,\Phi\,\bar\psi_L\psi_R + y\,\Phi\,\bar\eta_L\eta_R + \mathrm{h.c.}\big)
 - \lambda\,(\Phi^\dagger\Phi)^2 .
\end{align}
There is no tree mass term (scale invariance). The Yukawa is gauge invariant
($q_\Phi+q_{\psi_R}-q_{\psi_L}=0$). Along the flat radial direction the tree potential is
$V_{\rm tree}=(\lambda/4)\chi^4$.

\paragraph{Parameters.}
The independent parameters are the dark gauge coupling $\gD$, the dark Yukawa $y$, and the
CW vacuum expectation value $w=\langle\chi\rangle$ (GeV). Dependent quantities follow from CW:
\begin{equation}
\beta_\lambda=\frac{1}{8\pi^2}\Big(3\gD^4-8\,(y/\sqrt2)^4\Big),\quad
\mZp=\gD\,w,\quad
m_\chi=\frac{y\,w}{\sqrt2},\quad
m_{\rm scalon}=\sqrt{\beta_\lambda}\,w,
\end{equation}
with $\lambda$ fixed by CW minimisation. The critic-agent enforced the existence
condition $\beta_\lambda>0$ as the hard cut $y^4<1.5\,\gD^4$ (bosonic $\Zp$ contribution
must dominate the fermionic one), and corrected the vev window from the proposer's
$[100,10^6]$~GeV to $[0.05,10]$~GeV so that the GW peak can fall in the nHz PTA band
(literature benchmarks span $92.9$~MeV--$10$~GeV). The Higgs portal $\lambda_{\rm portal}$,
the kinetic mixing $\epsilon$, and the renormalisation scale $\mu$ are fixed
($\lambda_{\rm portal}=0$, $\epsilon=0$, $\mu=w$, no RG running, no Daisy). The branch-specific
feature is that all masses are proportional to $w$, mapping cleanly onto the
single-scalar semi-analytic FOPT backend.

\section{Pipeline status}
\label{sec:status}
Table~\ref{tab:status} summarises the status of every stage. The proposer model was
revised by the critic (vev window and $\beta_\lambda>0$ cut). The FOPT, PTA, constraint
and prior-audit stages all completed but carry \texttt{warning} status, which propagates
to the overall verdict.

\begin{table}[h]
\centering
\small
\caption{Pipeline stage status (from the corresponding handoff JSONs).}
\label{tab:status}
\begin{tabular}{lll}
\toprule
Stage / handoff & Status & Key note \\
\midrule
\texttt{handoff\_model} (critic) & ok & vev window and $\beta_\lambda>0$ cut applied \\
\texttt{handoff\_critique}       & ok & no red flags; non-novel (closest 2502.19478) \\
\texttt{handoff\_librarian\_preliminary} & ok & 5 papers; not novel \\
\texttt{handoff\_fopt}           & warning & $\betaH$ convention ($152.7$ vs $21.1$); $\alpha$ up to $10^{16}$ \\
\texttt{handoff\_pta}            & warning & $K_{sw}\to1$ saturation; no Bayes factor \\
\texttt{handoff\_constraints}    & warning & $\epsilon=0$ vs reheating; 5 missing calculations \\
\texttt{handoff\_prior\_audit}   & warning & 22 items; top severity $A01$ ($\epsilon$/reheating) \\
\bottomrule
\end{tabular}
\end{table}

\section{Preliminary literature search}
\label{sec:lit}
The librarian verdict is \emph{not novel: a close precedent exists}. The closest match,
arXiv:2502.19478 (Balan et al.), is essentially the same model including the Yukawa
dark fermion as sub-GeV dark matter, differing by its use of a kinetic-mixing portal.
The only distinctive ingredients of the target are the explicit mirror fermion $\eta$
for anomaly cancellation and the deliberate neglect of kinetic mixing for the FOPT.
Table~\ref{tab:lit} lists the verified papers; all five arXiv identifiers were
independently confirmed against arXiv/INSPIRE-HEP during this report.

\begin{table}[h]
\centering
\small
\caption{Verified preliminary-literature papers and classification.}
\label{tab:lit}
\begin{tabular}{p{1.7cm}p{6.4cm}p{4.3cm}c}
\toprule
arXiv & Title (abbrev.) & Relation to target & Score \\
\midrule
2502.19478 & Sub-GeV DM and nHz GW from a classically conformal dark sector & same model (adds kinetic mixing) & 10 \\
2501.11619 & Supercooled PTs in conformal dark sectors explain NANOGrav & structurally similar (scalar+gauge) & 9 \\
2501.15649 & Supercooled dark scalar PT explanation of NANOGrav & structurally similar (GeV vev) & 9 \\
2105.01007 & GW from a dark $U(1)_D$ PT in light of NANOGrav 12.5yr & same target, non-conformal & 7 \\
2306.17086 & NANOGrav signal from a dark conformal PT & background context (strong dynamics) & 5 \\
\bottomrule
\end{tabular}
\end{table}

\paragraph{Benchmark hints.}
The recommended anchor is arXiv:2502.19478 Point~A ($\gD=0.877$, $y=0.363$, $v=92.9$~MeV,
$T_p=8.36$~MeV, $\Treh=12.9$~MeV, $\alpha=4.99$, $\betaH=21.1$), which maps directly onto
$(\gD,y,w)$. Point~B ($\gD=0.79$, $y=0.401$, $v=343$~MeV, $\alpha=62.6$, $\betaH=13.3$) gives
stronger supercooling. arXiv:2501.15649 BP1--BP4 ($\gD\sim0.75\text{--}0.86$, $w=0.5\text{--}10$~GeV,
$\alpha\sim10^2\text{--}10^3$) bracket the GeV window; arXiv:2501.11619 ($\gD=0.59$, $\alpha\sim10^5$)
brackets the deep-supercooling end. The librarian flags non-uniform conventions
($\betaH$ vs $H_*R_*$; $T_p$ vs $T_*$ vs $\Treh$; CW-fixed vs free $\lambda$) that must be
reconciled before transferring numbers --- a caveat that materialises in the $\betaH$ mismatch below.

\section{FOPT results}
\label{sec:fopt}
The FOPT stage ran the semi-analytic backend (high-$T$ bounce action, Gaussian percolation,
$\mu=w$, no RG/Daisy; scalar and Goldstone contributions neglected, $g_b=0$) over $374$ points,
of which $238$ are \emph{viable} and $136$ failed/excluded. The mapping is
$\chi_0=w$, $g_b(\Zp)=\gD$ (dof 3), $g_f(\chi)=y/\sqrt2$ (dof 8).

The literature anchor (Point~A) is reproduced within $\sim10\text{--}15\%$ on $T_p$, $\Treh$ and
$\alpha$, but the backend reports $\betaH\simeq152.7$ versus the literature $21.1$ --- a factor
$\sim7$ \emph{convention mismatch} (severity-9 audit item $A02$). Deeply supercooled in-band points
reach $\alpha\sim10^{16}$ (vacuum domination), with $\betaH\simeq12\text{--}17$ and finite, tiny $T_p$.
Table~\ref{tab:fopt} lists the anchor and representative best points.
Validation flags (finite values, percolation, eternal-inflation, perturbativity) are all true.
The reported $f_{\rm peak}^{\rm est}$ is explicitly a template-independent order-of-magnitude redshifted
peak, \emph{not} the true spectral peak (which comes only from the PTA-stage template; audit $A21$).

\begin{table}[h]
\centering
\small
\caption{FOPT benchmarks (from \texttt{handoff\_fopt.json}); $w$ in GeV, temperatures in GeV.}
\label{tab:fopt}
\begin{tabular}{lcccccc}
\toprule
Point & $\gD$ & $y$ & $w$ & $\alpha$ & $\betaH$ & $\Treh$ \\
\midrule
lit anchor A & 0.877 & 0.363 & 0.0929 & $4.7$ & $152.7$ & $1.15\times10^{-2}$ \\
best\_1 & 0.50 & 0.20 & 0.121 & $7.5\times10^{16}$ & $12.3$ & $9.0\times10^{-3}$ \\
best\_3 & 0.50 & 0.00 & 0.707 & $7.4\times10^{15}$ & $12.4$ & $5.3\times10^{-2}$ \\
best\_7 & 0.54 & 0.233 & 0.484 & $3.6\times10^{12}$ & $16.2$ & $3.9\times10^{-2}$ \\
\bottomrule
\end{tabular}
\end{table}

\paragraph{Caveats.} The high-$T$ expansion is least reliable exactly in the deep-supercooling
regime where the model operates ($A11$); the Gaussian percolation approximation can mis-estimate
completion at the extreme-$\alpha$ tail ($A12$); the redshift uses SM $g_*(T)$ although with
$\epsilon=0$ the carrying bath is dark ($A18$). The GW-spectra overview figure is
\texttt{fopt\_gw\_spectra.png}.

\section{PTA results}
\label{sec:pta}
The PTA stage compared the model against the NG15 free spectrum ($14$ frequency bins) in two
modes: a benchmark violin scan and a full PTArcade/ceffyl Bayesian campaign.

\paragraph{Template.} The dissipative bulk-flow (\emph{dbf}) template was used (with a
bulk-flow \emph{bf} cross-check giving a near-identical peak), justified by the strongly
supercooled, vacuum-dominated regime with runaway walls $v_w\to1$; the higgsless sound-wave
template is invalid for $\alpha\gg1$. Because $\alpha\gg1$, the efficiency saturates
($K_{sw}\to1$): the amplitude is pinned at its strong-PT ceiling and is effectively
insensitive to the true (extreme) $\alpha$, so the fit constrains essentially $\betaH$ and
$\Treh$ rather than $\alpha$ (audit $A03$).

\paragraph{Benchmark comparison.} Of $945$ scanned points, the best score $14/14$ violin bins
was achieved by $29$ interior points. The four ranked best-fit points
(Table~\ref{tab:pta}) all score $14/14$, with peak frequencies $\sim10^{-8}$~Hz and
$\OmGW\sim$ few$\times10^{-9}$ at the peak. The benchmark spectrum overlaid on the
PTA violins is shown in Fig.~\ref{fig:bench}.

\begin{table}[h]
\centering
\small
\caption{PTA best-fit benchmark points (from \texttt{handoff\_pta.json}); $w$ in GeV.}
\label{tab:pta}
\begin{tabular}{llcccccc}
\toprule
Rank & $\gD$ & $y$ & $w$ & $\alpha$ & $\betaH$ & bins & $f_{\rm peak}$ [Hz] \\
\midrule
1 & 0.507 & 0.00 & 1.371 & $1.6\times10^{15}$ & 12.7 & 14/14 & $1.90\times10^{-8}$ \\
2 & 0.602 & 0.35 & 0.500 & $3.1\times10^{8}$  & 23.6 & 14/14 & $1.27\times10^{-8}$ \\
3 & 0.686 & 0.683 & 2.000 & $1.4\times10^{11}$ & 18.0 & 14/14 & $2.96\times10^{-8}$ \\
4 & 0.769 & 0.85 & 2.828 & $3.7\times10^{8}$  & 20.7 & 14/14 & $2.32\times10^{-8}$ \\
\bottomrule
\end{tabular}
\end{table}

\paragraph{PTArcade campaign.} The ceffyl campaign on NG15 ($10^5$ samples, acceptance $0.269$)
used the \emph{user-enlarged} priors $\gD\sim U(0.1,2.0)$, $y\sim U(0.0,2.0)$,
$\log_{10}(w/\mathrm{GeV})\sim U(-1.301,1.0)$, with the hard cut $y^4<1.5\,\gD^4$. The chains are
finite and the posteriors are peaked and \emph{interior}, with no prior-edge contact at $3\sigma$:
the enlarged upper edges ($\gD,y\to2.0$) carry no posterior mass. The Bayesian estimates are
\begin{equation}
\gD=0.588\pm0.060,\qquad y=0.377\pm0.222,\qquad \log_{10}(w/\mathrm{GeV})=-0.002\pm0.182\ (w\simeq1\ \mathrm{GeV}),
\end{equation}
with $68/95/99.7\%$ credible regions all interior. The MAP point
($\gD=0.535$, $y=0.294$, $\log_{10}w=0.097$, i.e. $w=1.25$~GeV) yields
$\alpha=1.8\times10^{15}$, $\betaH=13.0$, $\Treh=98.6$~MeV and scores $14/14$ bins.
\textbf{No Bayes factor / evidence was computed} (\texttt{evidence=null}); the ``fit'' is a
goodness-of-region and posterior-shape statement, not a model comparison against SMBHB or other
SGWB hypotheses (audit $A04$). The PTArcade spectrum (benchmark vs MAP, overlaid on violins) is in
Fig.~\ref{fig:ptarcade}, and the posterior distributions in Fig.~\ref{fig:posteriors}.

\begin{figure}[h]
\centering
\includegraphics[width=0.72\textwidth]{pta_benchmark_spectrum.pdf}
\caption{Estimated benchmark GW spectrum vs frequency, overlaid with the NG15 PTA violins.
Source: \texttt{pta\_benchmark\_spectrum.pdf}.}
\label{fig:bench}
\end{figure}

\begin{figure}[h]
\centering
\includegraphics[width=0.72\textwidth]{pta_ptarcade_spectrum.pdf}
\caption{PTArcade Bayesian MAP-point GW spectrum vs frequency, overlaid with the NG15 PTA violins.
Source: \texttt{pta\_ptarcade\_spectrum.pdf}.}
\label{fig:ptarcade}
\end{figure}

\begin{figure}[h]
\centering
\includegraphics[width=0.78\textwidth]{pta_posteriors.pdf}
\caption{Bayesian posterior distributions on $(\gD,y,\log_{10}w)$ from the PTArcade/ceffyl
campaign on NG15. Posteriors are peaked and interior; enlarged prior edges carry no mass.
Source: \texttt{pta\_posteriors.pdf}.}
\label{fig:posteriors}
\end{figure}

\section{Constraints}
\label{sec:constraints}
The constraint stage considered $12$ constraints (Table~\ref{tab:constr}): $1$ directly testable,
$2$ approximately testable, $3$ requiring recast, $5$ requiring a new backend, and $1$ not applicable.
The MAP point is \emph{provisionally allowed but unverified}. The decisive issue is that the FOPT/PTA
analyses fixed $\epsilon=0$, making the dark sector secluded: with strictly $\epsilon=0$ the latent
heat stays dark-only and the model would fail BBN/$N_{\rm eff}$. A nonzero kinetic mixing
$\epsilon\sim10^{-6}\text{--}10^{-3}$ is physically required, and that same window triggers the
dark-photon laboratory, beam-dump and SN1987A constraints --- none of which can be evaluated without
$\Zp$ decay widths and the portal energy-transfer calculation.

\begin{table}[h]
\centering
\small
\caption{Constraints (from \texttt{handoff\_constraints.json}). Class:
D=direct, A=approximate, R=recast, B=new backend, NA=not applicable.}
\label{tab:constr}
\begin{tabular}{clcp{5.6cm}c}
\toprule
ID & Constraint & Class & Brief & Source \\
\midrule
C1 & BBN reheating ($\Treh>$ few MeV) & A & SM bath must reheat above $\sim3$--4 MeV; needs portal transfer & 2502.19478 \\
C2 & Kinetic-mixing thermalisation & B & $\epsilon$ required to thermalise dark with SM & 2502.19478 \\
C3 & $\Delta N_{\rm eff}$ (dark radiation) & B & massless $\chi$ ($y{=}0$), scalon, $\Zp$ residue & 2502.19478 \\
C4 & $\Zp$ visible-decay limits & R & BaBar/KLOE/LHCb/NA48/NA64; excludes large $\epsilon$ & 2603.08430 \\
C5 & Beam-dump (E137,CHARM,$\nu$-Cal) & R & low-$\epsilon$ long-lived tail & 2603.08430 \\
C6 & SN1987A cooling & A & light $\mZp$, $\epsilon\sim10^{-10}$--$10^{-6}$ & 1611.03864 \\
C7 & $\chi$ relic / overclosure & B & $\Omega_\chi h^2\le0.12$; secluded annihilation & 2502.19478 \\
C8 & DM direct detection & B & $\sim\epsilon^2\gD^2$; depends on $f_\chi$ & 2603.08430 \\
C9 & Scalon decay / BBN injection & B & light scalon, loop-suppressed, possibly long-lived & 2502.19478 \\
C10 & Perturbativity / $\beta_\lambda>0$ & D & $\gD,y$ $O(1)$; cut enforced. \emph{Satisfied.} & internal \\
C11 & Percolation / no eternal inflation & D & checked upstream. \emph{Satisfied.} & internal (FOPT) \\
C12 & Higgs portal (invisible decay) & NA & $\lambda_{\rm portal}=0$ & internal \\
\bottomrule
\end{tabular}
\end{table}

Five missing calculations are flagged as required before any point can be declared allowed:
(i) portal energy transfer / SM reheating temperature; (ii) $\Zp$ decay widths/BR/lifetime;
(iii) $\Delta N_{\rm eff}$; (iv) $\chi$ relic abundance; (v) scalon decay/lifetime.

\section{Assumptions and limitations}
\label{sec:assumptions}
The prior-audit ranked $22$ items. Table~\ref{tab:assump} lists the highest-severity ones.
The two critical, decision-driving items are $A01$ (the $\epsilon=0$ vs reheating tension,
severity $10$, uncontrolled) and $A02$ (the $\betaH$ convention mismatch, severity $9$,
uncontrolled). Major limiting items include $A03$ (template $K_{sw}\to1$ saturation, so the
fit is $\alpha$-insensitive), $A04$ (no Bayes factor), $A10$ (renormalisation-scale/Daisy
systematic), $A22$ (reheating-viable $\epsilon$ may itself be excluded by lab/SN bounds),
and $A16$ (the nHz signal is window-conditional on the critic-chosen $w\in[0.05,10]$~GeV).

\begin{table}[h]
\centering
\small
\caption{Highest-severity assumptions/limitations (from \texttt{handoff\_prior\_audit.json}).
Ctrl = controlled.}
\label{tab:assump}
\begin{tabular}{clccl}
\toprule
ID & Description (brief) & Sev & Ctrl & Source \\
\midrule
A01 & $\epsilon=0$ secluded; reheating/BBN need $\epsilon\sim10^{-6}$--$10^{-3}$ & 10 & no & constraints \\
A02 & $\betaH$ convention mismatch ($152.7$ vs $21.1$) & 9 & no & fopt \\
A03 & dbf $K_{sw}\to1$ saturation; fit $\alpha$-insensitive & 8 & no & pta \\
A04 & no Bayes factor / evidence computed & 8 & no & pta \\
A22 & reheating-viable $\epsilon$ may be excluded by lab/SN & 7 & no & constraints \\
A05 & $\Delta N_{\rm eff}$ uncomputed (worst at $y{\to}0$) & 7 & no & constraints \\
A10 & $\mu=w$, no RG, no Daisy systematic on $T_p,\alpha,\beta$ & 7 & no & model \\
A11 & high-$T$ bounce action in deep supercooling & 6 & no & fopt \\
A13 & runaway walls $v_w\to1$ assumed & 6 & no & pta \\
A16 & nHz signal window-conditional on $w\in[0.05,10]$ GeV & 5 & yes & critique \\
A07 & enlarged priors ($\gD,y\le2.0$); posterior interior & 5 & yes & pta \\
A17 & non-novel; closest precedent 2502.19478 & 3 & yes & librarian \\
\bottomrule
\end{tabular}
\end{table}

\section{Inconsistencies and reference caveats}
\label{sec:inconsistencies}
One reference-content inconsistency was identified during citation validation:
\begin{itemize}
\item The PTA and prior-audit handoffs describe arXiv:2511.15687 as the source of the
``dbf bulk-flow'' GW template and as an ``expansion-corrected supercooled template''.
The arXiv identifier is \emph{valid and exists}, but the paper is actually
\emph{``Impact of cosmic expansion on gravitational wave spectra from strongly supercooled
first-order phase transitions''} by M.~Lewicki and V.~Vaskonen --- a paper on cosmic-expansion
corrections to supercooled-PT spectra, \emph{not} a dedicated dissipative-bulk-flow template
paper. We therefore cite arXiv:2511.15687 only for its actual subject (cosmic-expansion
corrections / cross-check context) and flag that the ``dbf template'' attribution in the
upstream handoffs is not corroborated by the cited paper's content. The numerical PTA results
are unaffected (the template was implemented in the backend code), but the bibliographic
attribution should be treated with caution.
\end{itemize}
No other numerical conflicts between handoffs were found; the FOPT $\betaH$ mismatch is an
internal-convention issue (reported as $A02$), not a cross-handoff data conflict.

\section{Conclusions}
\label{sec:conclusions}
\textbf{Can the model explain NANOGrav through a FOPT?}
\emph{For the signal fit: yes, conditionally.} Within the critic-chosen $w\in[0.05,10]$~GeV window,
the strongly supercooled FOPT produces a nHz SGWB that fits the NG15 free spectrum: the MAP point
and the four ranked benchmarks each cover $14/14$ violin bins, the Bayesian posteriors are peaked
and interior (no enlarged-prior-edge mass), and the preferred parameters
($\gD\simeq0.59$, $y\simeq0.38$, $w\simeq1$~GeV) are consistent with the closest literature
precedent (arXiv:2502.19478).

\emph{For cosmological viability: not established.} The result is conditional and partly
uncontrolled. (1) The whole chain fixed $\epsilon=0$; SM reheating and BBN require
$\epsilon\sim10^{-6}\text{--}10^{-3}$, and that window also triggers dark-photon lab/SN bounds that
were not evaluated ($A01$, $A22$, C1--C6). (2) A factor-$\sim7$ $\betaH$ convention mismatch
($A02$) and the saturated $K_{sw}\to1$ template ($A03$) mean the fit effectively constrains
$\betaH$ and $\Treh$ but not $\alpha$, and may be systematically biased. (3) No Bayes factor was
computed ($A04$), so ``explains'' is a goodness-of-region statement, not a model-comparison result
against SMBHB. (4) Five decisive cosmological observables remain uncomputed.

\emph{Robust conclusions:} the model is anomaly-free and self-consistent; couplings are perturbative
and $\beta_\lambda>0$ throughout (C10); the transition percolates without eternal inflation (C11);
the GW peak lands in the PTA band for the chosen $w$ window; and the PTA fit quality is high and
interior. \emph{Conditional/blocked:} cosmological viability, pending a kinetic-mixing portal,
reheating, $\Delta N_{\rm eff}$, relic-abundance, scalon-decay, and $\betaH$-reconciliation
calculations, plus an evidence/Bayes-factor determination.

\bibliographystyle{plain}
\begin{thebibliography}{9}
\bibitem{Balan2025} S.~Balan, T.~Bringmann, F.~Kahlhoefer, J.~Matuszak, C.~Tasillo,
\emph{Sub-GeV dark matter and nano-Hertz gravitational waves from a classically conformal dark sector},
JCAP \textbf{08} (2025) 062, arXiv:2502.19478.
\bibitem{Goncalves2025} J.~Gon\c{c}alves, D.~Marfatia, A.~P.~Morais, R.~Pasechnik,
\emph{Supercooled phase transitions in conformal dark sectors explain NANOGrav data},
Phys.\ Lett.\ B \textbf{869} (2025) 139829, arXiv:2501.11619.
\bibitem{Costa2025} F.~Costa, J.~Hoefken Zink, M.~Lucente, S.~Pascoli, S.~Rosauro-Alcaraz,
\emph{Supercooled Dark Scalar Phase Transitions explanation of NANOGrav data},
Phys.\ Lett.\ B \textbf{868} (2025) 139634, arXiv:2501.15649.
\bibitem{Borah2021} D.~Borah, A.~Dasgupta, S.~K.~Kang,
\emph{Gravitational waves from a dark $U(1)_D$ phase transition in the light of NANOGrav 12.5 yr data},
Phys.\ Rev.\ D \textbf{104} (2021) 063501, arXiv:2105.01007.
\bibitem{Fujikura2023} K.~Fujikura, S.~Girmohanta, Y.~Nakai, M.~Suzuki,
\emph{NANOGrav Signal from a Dark Conformal Phase Transition},
Phys.\ Lett.\ B (2023), arXiv:2306.17086.
\bibitem{Caputo2026} A.~Caputo, R.~Essig,
\emph{The Dark Photon: a 2026 Perspective}, arXiv:2603.08430.
\bibitem{Chang2017} J.~H.~Chang, R.~Essig, S.~D.~McDermott,
\emph{Revisiting Supernova 1987A Constraints on Dark Photons},
JHEP \textbf{01} (2017) 107, arXiv:1611.03864.
\bibitem{Lewicki2025} M.~Lewicki, V.~Vaskonen,
\emph{Impact of cosmic expansion on gravitational wave spectra from strongly supercooled
first-order phase transitions}, arXiv:2511.15687. (Cited for cosmic-expansion corrections;
see Sec.~\ref{sec:inconsistencies} for the upstream ``dbf template'' attribution caveat.)
\end{thebibliography}

\end{document}
